\documentclass[11pt,a4paper]{article}
\usepackage[utf8]{inputenc}
\usepackage{amsmath,amssymb}
\usepackage{geometry}
\geometry{margin=2.5cm}
\usepackage{listings}
\usepackage{xcolor}
\usepackage{hyperref}

\lstset{
  language=C++,
  basicstyle=\ttfamily\small,
  keywordstyle=\color{blue}\bfseries,
  commentstyle=\color{green!60!black},
  stringstyle=\color{red!70!black},
  numbers=left,
  numberstyle=\tiny\color{gray},
  numbersep=8pt,
  frame=single,
  breaklines=true,
  tabsize=4,
  showstringspaces=false
}

\title{IOI 1998 -- Picture}
\author{Solution}
\date{}

\begin{document}
\maketitle

\section{Problem Statement}

Given $n$ axis-aligned rectangles, compute the total perimeter of their union.
Only boundary segments separating the union's interior from the exterior are counted.

\medskip
\noindent\textbf{Constraints:} $1 \le n \le 5000$, coordinates in $[-10000, 10000]$.

\section{Solution Approach}

\subsection{Decomposition into Horizontal and Vertical Contributions}

The perimeter has horizontal segments (parallel to the $x$-axis) and vertical segments
(parallel to the $y$-axis). We compute each contribution separately using a sweep line.

\subsection{Horizontal Perimeter via Vertical Sweep}

\begin{enumerate}
  \item For each rectangle, create two events: at $y = y_1$ (bottom), \textbf{open} the
        interval $[x_1, x_2]$; at $y = y_2$ (top), \textbf{close} it.
  \item Sort events by $y$-coordinate. At the same $y$, process \textbf{open} events
        before \textbf{close} events (this avoids miscounting when rectangles share a
        horizontal edge).
  \item Maintain the total covered length along the $x$-axis. The horizontal perimeter
        contribution at each event $y$-coordinate is the absolute change in covered length:
        $|\Delta L|$.
\end{enumerate}

\subsection{Vertical Perimeter via Horizontal Sweep}

By symmetry, swap the roles of $x$ and $y$ and repeat the sweep.

\subsection{Implementation with Coordinate Compression}

For $n \le 5000$, an $O(n^2)$ approach suffices. We compress the $x$-coordinates
and maintain a coverage count array. At each event, we update the array and recompute
the total covered length.

\section{C++ Solution}

\begin{lstlisting}
#include <cstdio>
#include <cstring>
#include <cstdlib>
#include <algorithm>
#include <vector>
using namespace std;

struct Rect {
    int x1, y1, x2, y2;
};

struct Event {
    int y, x1, x2, type; // type: +1 = open, -1 = close
    bool operator<(const Event& o) const {
        if (y != o.y) return y < o.y;
        return type > o.type; // open (+1) before close (-1) at same y
    }
};

int n;
Rect rects[5001];

// Sweep in the y-direction, measuring horizontal perimeter
long long sweepHorizontal() {
    vector<Event> events;
    for (int i = 0; i < n; i++) {
        events.push_back({rects[i].y1, rects[i].x1, rects[i].x2, +1});
        events.push_back({rects[i].y2, rects[i].x1, rects[i].x2, -1});
    }
    sort(events.begin(), events.end());

    // Coordinate-compress x values
    vector<int> xs;
    for (int i = 0; i < n; i++) {
        xs.push_back(rects[i].x1);
        xs.push_back(rects[i].x2);
    }
    sort(xs.begin(), xs.end());
    xs.erase(unique(xs.begin(), xs.end()), xs.end());
    int M = xs.size();

    // cnt[j] = coverage count for interval [xs[j], xs[j+1])
    vector<int> cnt(M, 0);

    long long perimeter = 0;
    int prevCovered = 0;

    int i = 0;
    while (i < (int)events.size()) {
        int curY = events[i].y;
        // Process all events at this y-coordinate
        while (i < (int)events.size() && events[i].y == curY) {
            int x1 = events[i].x1, x2 = events[i].x2, t = events[i].type;
            for (int j = 0; j + 1 < M; j++) {
                if (xs[j] >= x1 && xs[j + 1] <= x2)
                    cnt[j] += t;
            }
            i++;
        }

        // Compute current covered length
        int covered = 0;
        for (int j = 0; j + 1 < M; j++)
            if (cnt[j] > 0)
                covered += xs[j + 1] - xs[j];

        perimeter += abs(covered - prevCovered);
        prevCovered = covered;
    }

    return perimeter;
}

// Sweep in the x-direction, measuring vertical perimeter
// Achieved by swapping x and y coordinates and reusing sweepHorizontal
long long sweepVertical() {
    Rect orig[5001];
    for (int i = 0; i < n; i++) orig[i] = rects[i];
    for (int i = 0; i < n; i++) {
        rects[i].x1 = orig[i].y1;
        rects[i].x2 = orig[i].y2;
        rects[i].y1 = orig[i].x1;
        rects[i].y2 = orig[i].x2;
    }
    long long result = sweepHorizontal();
    for (int i = 0; i < n; i++) rects[i] = orig[i]; // restore
    return result;
}

int main() {
    scanf("%d", &n);
    for (int i = 0; i < n; i++)
        scanf("%d %d %d %d", &rects[i].x1, &rects[i].y1,
              &rects[i].x2, &rects[i].y2);

    long long horiz = sweepHorizontal();
    long long vert = sweepVertical();

    printf("%lld\n", horiz + vert);
    return 0;
}
\end{lstlisting}

\section{Correctness}

The sweep line processes events sorted by $y$-coordinate. At each $y$, all intervals
opening or closing at that $y$ are applied to the coverage array. The change in total
covered length $|\Delta L|$ equals the total length of horizontal boundary segments
at that $y$-coordinate.

By processing opens before closes at the same $y$, we correctly handle rectangles
that share an edge: the coverage never drops to zero and back, avoiding double-counting.

The vertical contribution is computed identically after swapping coordinates.

\section{Complexity Analysis}

\begin{itemize}
  \item \textbf{Time complexity:} $O(n^2)$. There are $O(n)$ events, and at each event
        we update $O(n)$ intervals. Sorting is $O(n \log n)$.
        For an $O(n \log n)$ solution, a segment tree on compressed coordinates could
        replace the linear scan.
  \item \textbf{Space complexity:} $O(n)$ for events and coordinate arrays.
\end{itemize}

\end{document}
